\documentclass[../../main.tex]{subfiles}
\begin{document}

{\bf [解]}

題意より
\begin{align}
  a_1, a_2, b_1, b_2 \in \mathbb{Q} \label{eq:1}
\end{align}
である．このもとで考える．また，$c=|AB|$とおく．
三角形$OAB$は下図のようになる．

この時，$a,b,\cos\theta,\sin\theta$はピタゴラスの定理，余弦定理，三角形$OAB$の面積公式より
\begin{align}
  \begin{cases}
    a = \sqrt{a_1^2 + a_2^2} \\
    b = \sqrt{b_1^2 + b_2^2} \\
    c = |AB| = \sqrt{(a_1 - b_1)^2 + (a_2 - b_2)^2} \\
    \cos\theta = \dfrac{a^2 + b^2 - c^2}{2ab} \\
    \sin\theta = \dfrac{|a_1 b_2 - a_2 b_1|}{ab}
  \end{cases} \label{eq:2}
\end{align}
である．(i),(ii),(iii)の条件について調べる．

\subparagraph{(i) $ab\in \mathbb{Q}$の時}
$a^2, b^2, c^2 \in \mathbb{Q}$ と\cref{eq:2}から
\begin{align}
  \text{(i)} \to \text{(ii)}, \quad \text{(i)} \to \text{(iii)}
\end{align}
が成り立つ．

\subparagraph{(ii) $\cos\theta \in \mathbb{Q}$の時}

$\cos\theta \neq 0$ なら\cref{eq:2}において
\begin{align}
  ab = \dfrac{a^2 + b^2 - c^2}{2\cos\theta}
\end{align}
だから，$a^2,b^2,c^2$が有理数だから$ab$も有理数である．

次に$\cos\theta = 0$ なら$\sin\theta = \pm 1$ であり，
\cref{eq:2}において
\begin{align}
  ab = \dfrac{|a_1 b_2 - a_2 b_1|}{\sin\theta}
\end{align}
だから$ab$も有理数である．

以上より$\text{(ii)} \to \text{(i)}$ が成り立つ．

\subparagraph{(iii) $\sin\theta \in \mathbb{Q}$の時}

(ii)の時と同様に考える．
$\sin\theta \neq 0$ なら\cref{eq:2}において
\begin{align}
  ab = \dfrac{|a_1 b_2 - a_2 b_1|}{\sin\theta}
\end{align}
だから$ab$も有理数である．

次に$\sin\theta = 0$ なら$\cos\theta = \pm 1$ であり，
\cref{eq:2}において
\begin{align}
  ab = \dfrac{a^2 + b^2 - c^2}{2\cos\theta}
\end{align}
だから$a^2,b^2,c^2$が有理数だから$ab$も有理数である．

以上より$\text{(iii)} \to \text{(i)}$ が成り立つ．

\casesend

以上より$\text{(i)} \iff \text{(ii)}$，$\text{(i)} \iff \text{(iii)}$ だから，$\text{(ii)} \iff \text{(iii)}$ であり，題意は示された．

\end{document}
